\section{Holomorphic shell rigidity, normal jets, and an annular obstruction}
\label{sec:holomorphic-shell}
The finite-rank argument in this section removes the polynomial hypothesis
from the earlier shell-identification result. The domain must include the
centre of the measured spheres. We also count normal multiplicities and give
an explicit obstruction on a punctured domain. These are project derivations;
external priority and a proof-assistant formalization are not claimed.
They do not assert microscopic spectral universality.

\subsection{The analytic dilation determinant}
\begin{lemma}[Confluent dilation independence]\label{lem:dilation-jets}
Let $f_1,\ldots,f_r$ be linearly independent holomorphic functions on
$D_R=\{z:|z|<R\}$. Let $0<c_1<\cdots<c_s$, and let positive integers
$m_i$ satisfy $\sum_i m_i=r$. The determinant of the $r\times r$ matrix
\[
 \mathcal D(w)_{(i,k),j}=w^k f_j^{(k)}(c_iw),
 \qquad 0\le k<m_i,
\]
is not identically zero on $|w|<R/c_s$.
\end{lemma}
\begin{proof}
The finite-dimensional space spanned by the $f_j$ admits a basis with
strictly increasing orders of vanishing at zero. This follows by successively
eliminating the first nonzero Taylor coefficient; the process cannot leave a
nonzero function with every Taylor coefficient zero. A constant invertible
change of basis therefore gives
\[
 f_j(w)=a_jw^{n_j}+O(w^{n_j+1}),\qquad
 0\le n_1<\cdots<n_r,\quad a_j\ne0.
\]
Writing $(n)_k=n(n-1)\cdots(n-k+1)$, with $(n)_k=0$ for $k>n$, yields
\[
 \det\mathcal D(w)=\left(\prod_j a_j\right)w^{\sum_jn_j}
 \det\bigl[(n_j)_k c_i^{n_j-k}\bigr]_{(i,k),j}
 +O\bigl(w^{1+\sum_jn_j}\bigr).
\]
The constant determinant is nonzero. Otherwise a nonzero polynomial with at
most $r$ monomials would have positive zeros $c_i$ of total multiplicity at
least $r$. Such a polynomial has at most $r-1$ positive zeros counted with
multiplicity: divide by its lowest power, differentiate, and apply Rolle's
theorem inductively. For complex coefficients take a nonzero real part after
a common scalar rotation; all the specified zeros and derivatives persist.
Thus the displayed leading coefficient is nonzero. The constant change of
basis only multiplies the determinant by a nonzero constant.
\end{proof}

\subsection{A rank bound without a degree bound}
A holomorphic kernel on $D_{R_1}\times D_{R_2}$ has separated rank $r$ if $r$
is the least number of terms in an expression
\[
 K(z,w)=\sum_{j=1}^r a_j(z)b_j(w).
\]
In a minimal expression both lists of functions are linearly independent.
For $0<c<R_1R_2$, say that $K$ has multiplicity at least $m$ on $zw=c$ if
\[
 \partial_w^kK(z,c/z)=0\quad(0\le k<m)
\]
on the annulus $c/R_2<|z|<R_1$.

\begin{theorem}[Finite-rank contact bound]\label{thm:analytic-contact}
A nonzero holomorphic kernel of finite separated rank $r$ has total
multiplicity at most $r-1$ on distinct positive product curves $zw=c_i$.
\end{theorem}
\begin{proof}
If the total were at least $r$, reduce some multiplicities to make their sum
exactly $r$. On the common nonempty annulus
$c_s/R_2<|z|<R_1$, the column $(a_j(z))_j$ is annihilated by the matrix with
rows $(b_j^{(k)}(c_i/z))_j$. Multiplying row $(i,k)$ by $z^{-k}$ gives the
matrix in Lemma~\ref{lem:dilation-jets} with $w=1/z$. Its determinant is a
nonzero holomorphic function on $|w|<R_2/c_s$. It cannot vanish throughout the
open subannulus $1/R_1<|w|<R_2/c_s$. On a nonempty open set the matrix is
therefore invertible, so every $a_j$ vanishes there. The identity theorem
contradicts minimality of the separated representation.
\end{proof}

For a real-analytic diagonal $H(z)=K(z,\overline z)$, normal vanishing to
order $m$ on $|z|=\rho>0$ means
\[
 \left.\partial_t^k H(te^{i\theta})\right|_{t=\rho}=0
 \quad\hbox{for every $\theta$ and $0\le k<m$}.
\]
It is equivalent to the contact condition above with $c=\rho^2$.
To verify the direction we need, expand
$K(z,s/z)=\sum_{j\ge0}A_j(z)(s-c)^j$ locally near $s=c$.
The zeroth normal condition and the identity theorem give $A_0=0$.
Inductively, once $A_0,\ldots,A_{j-1}$ vanish, the $j$th normal derivative is
$j!(2\rho)^j A_j(\rho e^{i\theta})$. The same identity theorem gives
$A_j=0$. These identities continue to the full product-curve annulus.
The reverse direction follows by the same expansion. Consequently the sum
of positive zero-shell multiplicities of a nonzero rank-$r$ Hermitian
analytic kernel is at most $r-1$.

\subsection{Intensity identification in any complex dimension}
\begin{theorem}[Shell-and-jet identification]\label{thm:holomorphic-identify}
Let $P:B_R\subset\mathbb C^m\to\mathbb C^d$ and
$Q:B_R\to\mathbb C^e$ be holomorphic, with finite $d,e$. Let distinct radii
$0<\rho_i<R$ and positive integers $m_i$ satisfy $\sum_i m_i\ge d+e$.
Assume that the normal derivatives of
$\|P(tu)\|^2-\|Q(tu)\|^2$ of orders $0,\ldots,m_i-1$ vanish at $t=\rho_i$
for every unit vector $u$. Then
\[
 \langle P(z),P(w)\rangle=\langle Q(z),Q(w)\rangle
 \qquad(z,w\in B_R).
\]
There is a fixed linear isometry from the span of $P(B_R)$ onto the span of
$Q(B_R)$ taking $P(z)$ to $Q(z)$. In particular, if $d=e$, it extends to
$U\in U(d)$ and $Q=UP$ throughout the ball.
\end{theorem}
\begin{proof}
Restrict to a complex line $z=\zeta u$ and polarize the norm difference in
$\zeta,\overline\zeta$. This is a holomorphic kernel of separated rank at
most $d+e$. Theorem~\ref{thm:analytic-contact}, with the normal-jet equivalence,
forces the restricted diagonal to vanish identically. Every point of the
ball lies on such a line, so the norm difference is zero everywhere.
Expanding at zero, independence of the real-analytic monomials
$z^\alpha\overline z^\beta$ makes all polarized coefficients zero.
The polarized identity extends to the ball. It makes the map
$\sum_j t_jP(z_j)\mapsto\sum_jt_jQ(z_j)$ well-defined and norm preserving;
extending orthonormal bases proves the last assertion.
\end{proof}

In particular, $2d$ distinct spheres suffice for two $d$-output holomorphic
maps, with no polynomial, rationality, or finite-Taylor-support hypothesis.
Alternatively one sphere and its normal derivatives of orders $0$ through
$2d-1$ suffice. Mixed allocations of the same total multiplicity are allowed.
The actual separated rank of the norm-difference kernel can replace $d+e$.
All measurements here are functions on whole spheres, not finite angular
sample sets.

\begin{proposition}[Sharpness, including repeated shells]
In complex dimension one the count $2d$ cannot be decreased, even for
polynomials, and even if the $2d-1$ conditions are prescribed as arbitrary
positive shell multiplicities.
\end{proposition}
\begin{proof}
Let $\sum_i m_i=2d-1$ and write
\[
 B(x)=\prod_i(x-\rho_i^2)^{m_i}=\sum_{j=0}^{2d-1}\beta_jx^j.
\]
Its coefficients are nonzero and alternate in sign, with exactly $d$ of
each sign. Assign the components $\sqrt{\beta_j}z^j$ with $\beta_j>0$ to
$P$, and $\sqrt{-\beta_j}z^j$ with $\beta_j<0$ to $Q$. Then
$\|P(z)\|^2-\|Q(z)\|^2=B(|z|^2)$ has exactly the prescribed normal
multiplicities, but is not identically zero. No constant unitary relates the
maps. Multiplying both maps by $e^z$ gives nonpolynomial entire sharp examples.
\end{proof}

The same identification holds for rational maps with poles off the measured
spheres: multiply both maps by a common polynomial denominator nonzero there.
It also holds for meromorphic maps whenever such a holomorphic common
denominator is available. Multiplication by its squared modulus preserves
normal zero multiplicities. A finite-order pole at the centre can therefore
be cleared; an essential singularity cannot in general be ignored.

\subsection{Why the centre hypothesis matters}
The following explicit scalar construction has equal intensity on infinitely
many concentric circles without a constant phase relation. It is holomorphic
on $\mathbb C^*$, not at the centre.
For $q=1/4$, define the normally convergent product
\[
 \Theta(z;q)=\prod_{j=0}^{\infty}(1-q^jz)(1-q^{j+1}/z),\qquad z\ne0.
\]
Reindexing the product gives
$\Theta(qz;q)=-z^{-1}\Theta(z;q)$. Its zeros are exactly $q^{\mathbb Z}$.
Take $a=2,b=3,c=5,d=6/5$, so $ab=cd$, and set
\[
 \begin{split}
 P(z)&=\Theta(az;q)\Theta(bz;q)\Theta(c/z;q)\Theta(d/z;q),\\
 Q(z)&=P(1/z).
 \end{split}
\]
Both functions have real Laurent coefficients. The functional equation gives
\[
 P(qz)=q^{-2}z^{-4}P(z),\qquad
 Q(qz)=q^{-2}z^{-4}Q(z).
\]
Thus $R=P/Q$ is a meromorphic $q$-periodic function and
$R(1/z)=1/R(z)$. On $|z|=q^{n/2}$, $\overline z=q^n/z$, so
\[
 \overline{R(z)}=R(\overline z)=R(q^n/z)=1/R(z).
\]
Away from zeros and poles this proves $|R|=1$. Continuity of $|P|^2-|Q|^2$
extends the equality to the whole circle. The ratio is not constant: at
$z=1/3$, $P$ is zero whereas the four arguments defining $Q$ are
$5/3,2/5,6,9$, none in $q^{\mathbb Z}$. Hence $Q(1/3)\ne0$.
This supplies infinitely many equal-intensity shells for scalar functions on
a punctured domain. It does not contradict Theorem~\ref{thm:holomorphic-identify}.

\subsection{Relation to earlier work and verification scope}
The scalar two-circle uniqueness case is already contained in Perez's
Theorem 2.2 \cite{Perez2021}, which treats the Nevanlinna class and boundary
radial values; holomorphic functions on a larger disk are a special case.
We do not claim that scalar statement as new. Al Helal, Lebl and Nandi
\cite{ALN2024}, Theorems 1.4 and 3.5--3.15, study rational sphere maps and
prove degree-dependent classification results. Their norm-to-unitary step is
established prior mathematics, also proved above. The extension relative to
this repository is the confluent finite-rank analytic contact argument and
its vector-valued shell-and-jet consequences, not the unitary Gram argument,
the generalized Vandermonde principle, or the theta functional equation.
A literature search and a verified citation do not establish priority.

The accompanying verifier checks confluent determinants in two exact rational
implementations, leading Taylor jets of genuinely nonpolynomial functions,
prescribed sharp examples and negative controls. Theta-product evaluation is
separately labelled high-precision numerical evidence. None of those finite
checks replaces the analytic proof. The following section gives a separate
quantitative instability construction, rather than inferring a stability
estimate from exact uniqueness.
